\documentclass[11pt]{article}
\newcommand{\ListaTitulo}{Gabarito --- Lista 11}
% Preâmbulo compartilhado — MATD48, Listas de Exercícios 2026
% Usado por Lista{NN}.tex e Gabarito{NN}.tex via % Preâmbulo compartilhado — MATD48, Listas de Exercícios 2026
% Usado por Lista{NN}.tex e Gabarito{NN}.tex via % Preâmbulo compartilhado — MATD48, Listas de Exercícios 2026
% Usado por Lista{NN}.tex e Gabarito{NN}.tex via \input{preamble}

\usepackage[utf8]{inputenc}
\usepackage[T1]{fontenc}
\usepackage[brazil]{babel}
\usepackage{fullpage}
\usepackage{amsmath,amssymb}
\usepackage{enumitem}
\usepackage{xcolor}
\usepackage{fancyhdr}
\usepackage{titlesec}
\usepackage{listings}
\usepackage{hyperref}
\usepackage{booktabs}
\usepackage{graphicx}

\definecolor{matd48azul}{HTML}{0B4F6C}
\definecolor{matd48verde}{HTML}{2E8B57}

\titleformat{\section}{\normalfont\Large\bfseries\color{matd48azul}}{\thesection}{1em}{}
\titleformat{\subsection}{\normalfont\large\bfseries\color{matd48azul}}{\thesubsection}{1em}{}

\providecommand{\ListaTitulo}{Lista de Exercícios}

\setlength{\headheight}{14pt}
\pagestyle{fancy}
\fancyhf{}
\lhead{\small MATD48 --- Planejamento de Experimentos A}
\rhead{\small \ListaTitulo}
\cfoot{\thepage}
\renewcommand{\headrulewidth}{0.4pt}

\lstset{
  basicstyle=\ttfamily\small,
  breaklines=true,
  frame=single,
  columns=fullflexible,
  backgroundcolor=\color{gray!8},
  commentstyle=\color{matd48verde},
  keywordstyle=\color{matd48azul}\bfseries,
  language=R,
  extendedchars=true,
  literate=%
    {á}{{\'a}}1 {à}{{\`a}}1 {â}{{\^a}}1 {ã}{{\~a}}1
    {é}{{\'e}}1 {ê}{{\^e}}1 {í}{{\'i}}1 {ó}{{\'o}}1
    {ô}{{\^o}}1 {õ}{{\~o}}1 {ú}{{\'u}}1 {ç}{{\c c}}1
    {Á}{{\'A}}1 {À}{{\`A}}1 {Â}{{\^A}}1 {Ã}{{\~A}}1
    {É}{{\'E}}1 {Ê}{{\^E}}1 {Í}{{\'I}}1 {Ó}{{\'O}}1
    {Ô}{{\^O}}1 {Õ}{{\~O}}1 {Ú}{{\'U}}1 {Ç}{{\c C}}1
}

\hypersetup{colorlinks=true, linkcolor=matd48azul, urlcolor=matd48azul, citecolor=matd48azul}

\newcommand{\cabecalho}[2]{%
  \begin{center}
    {\Large\bfseries\color{matd48azul} #1}\\[0.3em]
    {\large #2}\\[0.3em]
    \rule{\linewidth}{0.6pt}
  \end{center}
  \vspace{0.5em}
}

\newenvironment{questao}[1]{\vspace{1em}\noindent\textbf{Questão #1.}\hspace{0.4em}}{\par}
\newenvironment{solucao}{\vspace{0.3em}\noindent\textit{Solução.}\hspace{0.4em}\color{black!85}}{\par\vspace{0.5em}}


\usepackage[utf8]{inputenc}
\usepackage[T1]{fontenc}
\usepackage[brazil]{babel}
\usepackage{fullpage}
\usepackage{amsmath,amssymb}
\usepackage{enumitem}
\usepackage{xcolor}
\usepackage{fancyhdr}
\usepackage{titlesec}
\usepackage{listings}
\usepackage{hyperref}
\usepackage{booktabs}
\usepackage{graphicx}

\definecolor{matd48azul}{HTML}{0B4F6C}
\definecolor{matd48verde}{HTML}{2E8B57}

\titleformat{\section}{\normalfont\Large\bfseries\color{matd48azul}}{\thesection}{1em}{}
\titleformat{\subsection}{\normalfont\large\bfseries\color{matd48azul}}{\thesubsection}{1em}{}

\providecommand{\ListaTitulo}{Lista de Exercícios}

\setlength{\headheight}{14pt}
\pagestyle{fancy}
\fancyhf{}
\lhead{\small MATD48 --- Planejamento de Experimentos A}
\rhead{\small \ListaTitulo}
\cfoot{\thepage}
\renewcommand{\headrulewidth}{0.4pt}

\lstset{
  basicstyle=\ttfamily\small,
  breaklines=true,
  frame=single,
  columns=fullflexible,
  backgroundcolor=\color{gray!8},
  commentstyle=\color{matd48verde},
  keywordstyle=\color{matd48azul}\bfseries,
  language=R,
  extendedchars=true,
  literate=%
    {á}{{\'a}}1 {à}{{\`a}}1 {â}{{\^a}}1 {ã}{{\~a}}1
    {é}{{\'e}}1 {ê}{{\^e}}1 {í}{{\'i}}1 {ó}{{\'o}}1
    {ô}{{\^o}}1 {õ}{{\~o}}1 {ú}{{\'u}}1 {ç}{{\c c}}1
    {Á}{{\'A}}1 {À}{{\`A}}1 {Â}{{\^A}}1 {Ã}{{\~A}}1
    {É}{{\'E}}1 {Ê}{{\^E}}1 {Í}{{\'I}}1 {Ó}{{\'O}}1
    {Ô}{{\^O}}1 {Õ}{{\~O}}1 {Ú}{{\'U}}1 {Ç}{{\c C}}1
}

\hypersetup{colorlinks=true, linkcolor=matd48azul, urlcolor=matd48azul, citecolor=matd48azul}

\newcommand{\cabecalho}[2]{%
  \begin{center}
    {\Large\bfseries\color{matd48azul} #1}\\[0.3em]
    {\large #2}\\[0.3em]
    \rule{\linewidth}{0.6pt}
  \end{center}
  \vspace{0.5em}
}

\newenvironment{questao}[1]{\vspace{1em}\noindent\textbf{Questão #1.}\hspace{0.4em}}{\par}
\newenvironment{solucao}{\vspace{0.3em}\noindent\textit{Solução.}\hspace{0.4em}\color{black!85}}{\par\vspace{0.5em}}


\usepackage[utf8]{inputenc}
\usepackage[T1]{fontenc}
\usepackage[brazil]{babel}
\usepackage{fullpage}
\usepackage{amsmath,amssymb}
\usepackage{enumitem}
\usepackage{xcolor}
\usepackage{fancyhdr}
\usepackage{titlesec}
\usepackage{listings}
\usepackage{hyperref}
\usepackage{booktabs}
\usepackage{graphicx}

\definecolor{matd48azul}{HTML}{0B4F6C}
\definecolor{matd48verde}{HTML}{2E8B57}

\titleformat{\section}{\normalfont\Large\bfseries\color{matd48azul}}{\thesection}{1em}{}
\titleformat{\subsection}{\normalfont\large\bfseries\color{matd48azul}}{\thesubsection}{1em}{}

\providecommand{\ListaTitulo}{Lista de Exercícios}

\setlength{\headheight}{14pt}
\pagestyle{fancy}
\fancyhf{}
\lhead{\small MATD48 --- Planejamento de Experimentos A}
\rhead{\small \ListaTitulo}
\cfoot{\thepage}
\renewcommand{\headrulewidth}{0.4pt}

\lstset{
  basicstyle=\ttfamily\small,
  breaklines=true,
  frame=single,
  columns=fullflexible,
  backgroundcolor=\color{gray!8},
  commentstyle=\color{matd48verde},
  keywordstyle=\color{matd48azul}\bfseries,
  language=R,
  extendedchars=true,
  literate=%
    {á}{{\'a}}1 {à}{{\`a}}1 {â}{{\^a}}1 {ã}{{\~a}}1
    {é}{{\'e}}1 {ê}{{\^e}}1 {í}{{\'i}}1 {ó}{{\'o}}1
    {ô}{{\^o}}1 {õ}{{\~o}}1 {ú}{{\'u}}1 {ç}{{\c c}}1
    {Á}{{\'A}}1 {À}{{\`A}}1 {Â}{{\^A}}1 {Ã}{{\~A}}1
    {É}{{\'E}}1 {Ê}{{\^E}}1 {Í}{{\'I}}1 {Ó}{{\'O}}1
    {Ô}{{\^O}}1 {Õ}{{\~O}}1 {Ú}{{\'U}}1 {Ç}{{\c C}}1
}

\hypersetup{colorlinks=true, linkcolor=matd48azul, urlcolor=matd48azul, citecolor=matd48azul}

\newcommand{\cabecalho}[2]{%
  \begin{center}
    {\Large\bfseries\color{matd48azul} #1}\\[0.3em]
    {\large #2}\\[0.3em]
    \rule{\linewidth}{0.6pt}
  \end{center}
  \vspace{0.5em}
}

\newenvironment{questao}[1]{\vspace{1em}\noindent\textbf{Questão #1.}\hspace{0.4em}}{\par}
\newenvironment{solucao}{\vspace{0.3em}\noindent\textit{Solução.}\hspace{0.4em}\color{black!85}}{\par\vspace{0.5em}}


\begin{document}

\cabecalho{Gabarito --- Lista de Exercícios 11}{Quadrado latino e quadrado greco-latino (Capítulo 4 / Aula 11)}

\begin{questao}{1} \textbf{(Agricultura --- ANOVA de um quadrado latino)}
\end{questao}
\begin{solucao}
\begin{enumerate}[label=(\alph*)]
  \item $SQ_{Erro} = SQ_{Total} - SQ_{Linha} - SQ_{Coluna} - SQ_{Variedade}
    = 300-40-35-150 = \textbf{75}$. Graus de liberdade: $gl_{Linha}=gl_{Coluna}=gl_{Var}=n-1=4$,
    $gl_{Erro}=(n-1)(n-2)=(4)(3)=12$, $gl_{Total}=n^2-1=24$ (confira: $4+4+4+12=24$
    \checkmark). Quadrados médios: $QM_{Var}=150/4=37{,}5$; $QM_{Erro}=75/12=6{,}25$.
  \item $F_0 = QM_{Var}/QM_{Erro} = 37{,}5/6{,}25 = \textbf{6{,}00}$.
  \item Como $F_0=6{,}00 > F_{0{,}05;4;12}=3{,}26$, \textbf{rejeita-se $H_0$}: há evidência de
    diferença entre as variedades de soja.
\end{enumerate}
\end{solucao}

\begin{questao}{2} \textbf{(Dedução --- graus de liberdade do erro em um quadrado latino replicado)}
\end{questao}
\begin{solucao}
\begin{enumerate}[label=(\alph*)]
  \item Cada uma das $p$ réplicas tem $n^2$ observações (uma por célula do quadrado $n\times n$),
    logo o total é $pn^2$ observações, e $gl_{Total}=pn^2-1$.
  \item Réplica: $p-1$ graus de liberdade. Linhas dentro de réplica: cada réplica contribui
    $n-1$ graus de liberdade (linhas são específicas de cada réplica, não compartilhadas), num
    total de $p(n-1)$. Colunas dentro de réplica: pelo mesmo argumento, $p(n-1)$. Tratamento:
    como o tratamento é comum a todas as réplicas, apenas $n-1$ graus de liberdade.
  \item Somando os termos do item (b):
    $$
    (p-1) + p(n-1) + p(n-1) + (n-1) = (p-1) + 2p(n-1) + (n-1).
    $$
    Subtraindo do total do item (a):
    $$
    gl_{Erro} = \big[pn^2-1\big] - \big[(p-1)+2p(n-1)+(n-1)\big]
              = pn^2 - p - 2p(n-1) - (n-1).
    $$
    Fatorando $p$ dos dois primeiros termos que o contêm:
    $$
    gl_{Erro} = p\big[n^2-1-2(n-1)\big] - (n-1) = p(n-1)(n+1-2) - (n-1) = p(n-1)^2 - (n-1),
    $$
    usando $n^2-1=(n-1)(n+1)$. Colocando $(n-1)$ em evidência:
    $$
    gl_{Erro} = (n-1)\big[p(n-1)-1\big].
    $$
  \item Para $n=4$, $p=2$: $gl_{Erro} = (3)\big[2(3)-1\big] = 3\times5 = \textbf{15}$ — o mesmo
    valor obtido pelo ajuste em R na Aula 11.
\end{enumerate}
\end{solucao}

\begin{questao}{3} \textbf{(Verificação --- validade de um quadrado latino)}
\end{questao}
\begin{solucao}
\begin{enumerate}[label=(\alph*)]
  \item As quatro linhas são, cada uma, uma permutação de \{A,B,C,D\} (nenhuma linha repete
    letra) — a condição de linha está satisfeita. Já as colunas: a coluna 3 tem
    (C, D, A, A) — as letras \textbf{A aparece duas vezes} (nas linhas 3 e 4), e a letra B nunca
    aparece nessa coluna; a coluna 4 tem (D, C, B, B) — a letra \textbf{B aparece duas vezes}
    (linhas 3 e 4), e a letra A nunca aparece nessa coluna. \textbf{O arranjo não é um quadrado
    latino válido}: a violação está exatamente nas linhas 3 e 4, que coincidem nas duas últimas
    colunas ("$\dots$A B").
  \item Basta trocar as duas últimas entradas da linha 4, de "A B" para "B A": a linha 4 passa de
    $D\ C\ A\ B$ para $D\ C\ B\ A$. O quadrado corrigido é
    $$
    \begin{array}{cccc}
    A & B & C & D \\
    B & A & D & C \\
    C & D & A & B \\
    D & C & B & A \\
    \end{array}
    $$
    Confira: toda linha é uma permutação de \{A,B,C,D\}, e as colunas agora são
    (A,B,C,D), (B,A,D,C), (C,D,A,B), (D,C,B,A) — todas também permutações completas, sem
    repetição. O quadrado corrigido é válido.
\end{enumerate}
\end{solucao}

\begin{questao}{4} \textbf{(Quadrado greco-latino --- existência e graus de liberdade)}
\end{questao}
\begin{solucao}
\begin{enumerate}[label=(\alph*)]
  \item Para $n=5$: $(n-1)(n-3)=(4)(2)=\textbf{8}$. Para $n=7$: $(n-1)(n-3)=(6)(4)=\textbf{24}$.
  \item Não existe quadrado greco-latino de ordem 6 porque não existe um par de quadrados latinos
    \textbf{ortogonais} de ordem 6 — este é o célebre resultado do "problema dos 36 oficiais" de
    Euler (conjecturado por Euler em 1782 e demonstrado rigorosamente por Tarry em 1900): é
    impossível arranjar 36 oficiais de 6 regimentos e 6 patentes em um quadrado $6\times6$ de modo
    que cada linha e coluna contenha exatamente um oficial de cada regimento e de cada patente.
  \item Um "quadrado hiper-greco-latino" com 3 conjuntos de letras sobrepostos (linha, coluna e
    3 fatores de tratamento, cada um com $n-1$ graus de liberdade, totalizando 5 fontes) deixaria
    $gl_{Erro} = n^2-1-5(n-1)$. Para $n=4$: $gl_{Erro}=15-5(3)=15-15=\textbf{0}$ graus de
    liberdade para o erro — \textbf{inviável na prática}: não sobra nenhuma informação para
    estimar a variabilidade residual, tornando impossível qualquer teste de hipótese. Esse desenho
    só se torna viável para $n$ maior (é necessário $n^2-1 > 5(n-1)$, ou seja $n>5$), e ainda
    exige a existência de 3 quadrados latinos mutuamente ortogonais de ordem $n$, o que nem sempre
    é garantido.
\end{enumerate}
\end{solucao}

\begin{questao}{5} \textbf{(Psicologia --- quadrado latino para contrabalancear ordem)}
\end{questao}
\begin{solucao}
\begin{enumerate}[label=(\alph*)]
  \item \textbf{Linha} = participante (4 participantes); \textbf{coluna} = posição de
    apresentação (1ª, 2ª, 3ª, 4ª sessão); \textbf{tratamento} = tipo de tarefa (A, B, C, D). O
    arranjo garante que cada participante realiza as 4 tarefas (uma vez cada) e que cada tarefa
    aparece exatamente uma vez em cada posição de apresentação ao longo dos 4 participantes —
    contrabalanceando tanto diferenças entre participantes quanto efeitos de ordem.
  \item Um DBCA com bloco = participante apenas controlaria as diferenças individuais entre
    participantes, mas \textbf{não controlaria o efeito de posição/ordem}: se a tarefa B fosse
    sistematicamente apresentada em 3º ou 4º lugar para a maioria dos participantes (por
    conveniência do protocolo, por exemplo), um eventual efeito de fadiga ficaria confundido com
    o efeito da tarefa B, sem que o DBCA tivesse como separá-los. O quadrado latino resolve isso
    tratando a posição como uma segunda fonte de bloqueio, cruzada com participante.
  \item A suposição de ausência de interação ordem$\times$tarefa falharia, por exemplo, se a
    tarefa D exigisse memória de trabalho intensa e fosse desproporcionalmente prejudicada por
    fadiga quando apresentada na 4ª posição (mas não nas demais tarefas, mais automáticas ou menos
    exigentes) — nesse caso, o efeito de "ser a tarefa D" dependeria de \emph{quando} ela foi
    apresentada, violando a aditividade que o modelo do quadrado latino assume.
\end{enumerate}
\end{solucao}

\begin{questao}{6} \textbf{(Aleatorização de Yates e análise de poder)}
\end{questao}
\begin{solucao}
\begin{enumerate}[label=(\alph*)]
  \item Reutilizar sempre o mesmo quadrado fixo torna a atribuição de tratamentos a células
    \textbf{determinística}, não aleatória: se, por coincidência, a diagonal principal (onde A
    sempre cai) se alinhar com alguma característica não modelada do processo (por exemplo, todas
    as máquinas na diagonal serem as mais novas, ou todos os horários na diagonal serem o turno da
    manhã), o efeito do tratamento A ficaria \textbf{confundido} com essa característica em
    \emph{todo} experimento que reusasse esse mesmo quadrado — exatamente o tipo de falha de
    delineamento estudado na Aula 09. Os quatro passos de Yates (1933) evitam isso: (1) partir de
    um quadrado latino reduzido; (2) permutar aleatoriamente as \textbf{linhas}; (3) permutar
    aleatoriamente as \textbf{colunas}; (4) sortear aleatoriamente quais \textbf{tratamentos
    reais} correspondem a cada símbolo $1,\dots,n$. O resultado é que cada execução do
    procedimento produz um arranjo diferente, todos igualmente prováveis, e é esse espaço de
    arranjos possíveis — não um único quadrado fixo — que sustenta a distribuição de referência
    sob a qual o teste $F$ é válido (o mesmo argumento de Neyman-Rubin da Aula 09, agora aplicado a
    uma aleatorização restrita simultaneamente por linha e por coluna).
  \item São necessárias $p=\textbf{5}$ réplicas para cruzar poder $\ge0{,}80$ (a tabela salta de
    $0{,}716$ com $p=4$ para $0{,}828$ com $p=5$). O trade-off: com $n=4$, cada réplica usa
    $n^2=16$ unidades experimentais, de modo que $p=5$ réplicas exigem $N=5\times16=\textbf{80}$
    unidades no total, contra apenas 16 unidades (e poder de só $14{,}8\%$) com um único quadrado —
    um aumento de 5 vezes no custo experimental para elevar o poder de detecção de um efeito
    grande de $14{,}8\%$ para $82{,}8\%$. Na prática, essa é uma decisão de custo-benefício: se
    recursos são escassos, o pesquisador pode aceitar um poder menor, buscar reduzir $\sigma^2$ por
    um controle experimental mais rígido (em vez de mais réplicas), ou reconsiderar se $f=0{,}4$ é
    de fato o menor efeito cientificamente relevante a detectar.
\end{enumerate}
\end{solucao}

\end{document}
